\section{Introduction}
\label{sec:intro}

A coding agent is a loop around a generative model with a handful of tools.
The design notes on Jev~\cite{jevnote} argue that the leverage in such a system
lies less in the loop than in what the loop feeds the model each time it turns,
and that underneath every feature of a capable harness sits a small decision
asked at high frequency. Should this chunk of context be visible to the next
turn, and in how much detail? Can this subtask leave the frontier model? Which
of several hundred tools fits the stated intent? May this command run? A
session asks such questions thousands of times.

Most agents answer them by asking a generative model and reading its prose, or
by a fixed rule. The generative answer costs a generation, has to be parsed,
and carries no probability that a caller could compare with a threshold. The
fixed rule is cheap and auditable but blind to everything its author did not
foresee. A \emph{typed decision} sits between the two. The caller states the
question, enumerates the admissible answers with a description of each, and
passes the state the answer should depend on. The model returns a probability
distribution over exactly those answers. The output can be validated against
its type, compared with a threshold, logged with the state it was computed
from, and replayed.

\paragraph{Three planes.}
Hanzo separates an agentic system into three planes with different kinds of
authority.
\begin{itemize}[leftmargin=1.5em]
  \item The \emph{generative plane} writes: code, plans, explanations, drafts.
    In Hanzo this is the Zen model family and the models it routes to.
  \item The \emph{decision plane} classifies: it answers typed questions about
    explicit state with calibrated distributions. This is Kai.
  \item The \emph{policy plane} decides what is permitted. It is deterministic,
    written by people, and has the last word.
\end{itemize}
The decision plane is advisory to the policy plane by construction
(Section~\ref{sec:authority}): a model's verdict can make an outcome stricter
than the policy's and never looser. The generative plane is advisory to both
where a decision governs it, such as which context it sees or whether its
command runs.

\paragraph{The decision as a program.}
Most AI systems treat a decision as text: evidence goes into a prompt, a
recommendation comes out as prose, and what produced it is not recorded in any
form a machine can check (Figure~\ref{fig:program}, left). Hanzo treats the
decision as a program. Generative models help build the program and reason
inside it. Typed decision models make the bounded judgments within it.
Deterministic code does the arithmetic. Policy and people hold authority. Every
run produces a reproducible record, a Decision Package, that can be replayed,
compared with another run, and refreshed when its evidence changes
(Figure~\ref{fig:program}, right). A typed decision is the primitive this
requires, between ordinary deterministic code and free-form generation. It is
more than a cheap classifier: it is the step that gives a judgment a type, a
probability and a record.

% Decision as text against decision as a program.
\begin{figure}[t]
\centering
\begin{tikzpicture}[
  font=\footnotesize,
  box/.style={draw, rounded corners=2pt, align=center, minimum height=1.5em,
              inner sep=3pt, fill=white},
  src/.style={box, fill=gray!8, minimum width=6.2em},
  exec/.style={box, fill=blue!6, minimum width=5.4em},
  auth/.style={box, fill=orange!10, minimum width=10em},
  arr/.style={-{Stealth[length=4pt]}, thin},
  node distance=0.55em and 1.2em]

% Left: the decision as text.
\node[font=\footnotesize\bfseries] (lt) {Decision as text};
\node[src, below=0.8em of lt] (le) {Evidence};
\node[box, below=of le] (lp) {One large prompt};
\node[box, below=of lp] (ll) {LLM};
\node[box, below=of ll] (lr) {Recommendation\\(prose)};
\draw[arr] (le) -- (lp);
\draw[arr] (lp) -- (ll);
\draw[arr] (ll) -- (lr);
\node[below=0.3em of lr, text=gray, align=center] {no typed record of\\what produced it};

% Right: the decision as a program.
\node[font=\footnotesize\bfseries, right=17em of lt] (rt) {Decision as a program};
\node[src, minimum width=4.6em, anchor=north east] (e2) at ([xshift=-0.2em,yshift=-0.8em]rt.south) {SysML / graph};
\node[src, minimum width=4.6em, anchor=north west] (e3) at ([xshift=0.2em,yshift=-0.8em]rt.south) {Constraints};
\node[src, minimum width=4.6em, left=0.4em of e2] (e1) {Evidence};
\node[src, minimum width=4.6em, right=0.4em of e3] (e4) {Assumptions};
\coordinate (srow) at (e2.south -| rt);
\node[box, fill=green!7, minimum width=12em, below=1.8em of srow] (dp)
  {\textbf{Decision Program}\\\scriptsize name@version, questions, thresholds};
\foreach \e in {e1,e2,e3,e4} \draw[arr] (\e.south) -- (dp.north);

\node[exec, below=1.4em of dp] (zen) {Zen\\\scriptsize reasoning};
\node[exec, left=0.7em of zen] (kai) {Kai\\\scriptsize typed judgment};
\node[exec, right=0.7em of zen] (det) {Deterministic\\\scriptsize code, solvers};
\draw[arr] (dp) -- (kai.north);
\draw[arr] (dp) -- (zen.north);
\draw[arr] (dp) -- (det.north);

\node[auth, below=1.4em of zen] (pol) {Policy\\\scriptsize deterministic authority};
\draw[arr] (kai.south) -- (pol);
\draw[arr] (zen.south) -- (pol);
\draw[arr] (det.south) -- (pol);
\node[auth, below=of pol] (hum) {Human approval};
\node[box, fill=green!7, minimum width=12em, below=of hum] (pkg)
  {\textbf{Decision Package}\\\scriptsize answers, versions, evidence, overrides};
\draw[arr] (pol) -- (hum);
\draw[arr] (hum) -- (pkg);
\draw[arr] (pkg.east) -- ++(2.6em,0) |- node[pos=0.25, right, align=left]
  {\scriptsize replay\\\scriptsize diff\\\scriptsize refresh} (dp.east);
\end{tikzpicture}
\caption{Left: the usual shape, in which one generative call turns evidence
into a recommendation and nothing typed records how. Right: Hanzo Decision.
A versioned program DAG states the questions; Kai answers the bounded
judgments, Zen the open reasoning, deterministic code and solvers the
arithmetic; policy and a person hold authority; every run produces a hashed
package that can be replayed, diffed and refreshed.}
\label{fig:program}
\end{figure}


The layers are separable, and each builds on the one below it. Jev put typed
questions inside software. Kai answers them natively, as open weights on
Hanzo's own runtime. Hanzo Decision adds versioned programs, policy, evidence,
provenance, refresh and replay. Hanzo's cloud deploys it as infrastructure.

\paragraph{The system.}
Hanzo Decision is the capability; \op{POST /v1/decisions} is its only public
interface; Kai is the model family behind it. The interface's wire format is
OpenRouter's Decisions API~\cite{openrouterdecisions}, so a client written for
TypeSafe's Jev~\cite{typesafejev} works unchanged, and the \op{model} field
chooses the backend: a \op{kai-*} name is answered natively and a
\op{typesafe/jev-*} name is forwarded to OpenRouter. The Kai checkpoints are
the three open Laya checkpoints~\cite{laya}, republished under
\texttt{hanzoai/} at pinned revisions with byte-identical weights, and a
native runtime for them is being written in Rust on Hanzo's tensor library so
that one binary can serve them on a CPU without Python.

\paragraph{Contributions.}
Each item is marked with its status: \meas{} for a figure from a run we made,
\reported{} for a figure another party published, \shipped{} for code or
weights that exist at the revision named, and \proposed{} for a design that is
not yet on a main branch.
\begin{enumerate}[leftmargin=1.5em]
  \item A definition of typed questions and of the uncertainty an answer
    reports, and the observation that two existing systems report different
    quantities under the name \op{confidence}, so a threshold tuned on one
    does not transfer to the other (Section~\ref{sec:decisions}).
  \item The Kai architecture as ported to Rust: the input layout, the typed
    head, temperature by question type and option count, and the act head,
    together with the provenance of each checkpoint
    (Section~\ref{sec:model}). The weights are \shipped; the runtime is
    \proposed.
  \item Policy precedence: a verdict lattice, abstention to \ask, three
    execution modes, and a proof that enforcing a model under the join can
    only add friction, which makes promotion a question answerable from shadow
    records (Section~\ref{sec:authority}). \proposed
  \item Decision Programs, the seven programs an agent loop needs, the decision
    record, and the model-agnostic wire (Section~\ref{sec:programs}).
    \proposed
  \item Decision-as-a-program for studies: the DecisionPackage, closed forms
    for the smallest weight change and the smallest answer change that flip a
    recommendation, and refresh that recomputes only answers whose evidence
    changed (Section~\ref{sec:studies}). \proposed
  \item Measurements of the three bootstrap checkpoints on nine recorded cases,
    and the protocol by which the rest of the system will be measured
    (Section~\ref{sec:evaluation}). \meas
\end{enumerate}

\paragraph{What this paper does not claim.}
It reports no accuracy of Kai on Hanzo's traffic, because no Kai checkpoint has
yet been trained on it. It reports no latency for the Rust runtime, because
that runtime has not yet been measured against the reference. It claims no cost
saving from moving decisions off the generative plane; Section~\ref{sec:eval-plan}
says how that will be measured. Throughput and accuracy figures for the
upstream checkpoints are quoted from their authors and marked as such.
